We studied when a diffusion language sampler should commit a position to the KV cache.
In DiffusionGemma, freezing tokens upon acceptance degrades performance at higher entropy bounds, while block-level commitment spends much of each window on positions that have already converged.
RES commits a sufficiently long accepted prefix and keeps positions farther ahead open to revision.
It needs no training, and across seven benchmarks it increases tokens per forward pass from 19.3 to 25.5, while macro accuracy goes from 80.7\% to 78.8\%.
The block sampler remains better at low throughput.

\paragraph{Limitations}
Tokens per forward pass counts forward passes, not wall-clock time, although each RES pass processes the same number of positions as the baseline.
All experiments use one model, DiffusionGemma.

\paragraph{Future Work}
The largest cost of RES is on instruction following, where IFEval accuracy falls by 10.7\% (absolute).
One possible cause is that RES creates a left-to-right noise pattern that differs from the noise seen in training.
Training with a rolling noise pattern, as in Clock Diffusion \citep{schiff2026clock}, is a natural way to test and remove this cost.
